% v4.8.0 release layers
% Appended layers v4.7.0-r85 to r94 and v4.8.0-r01, r02, carried verbatim from the full-monolith working source
% Quantum_Sphaera_Companion_v4_8_0_source.tex (SHA256 a1b988d3c9e3642f93e201bcb56d6ff03020ae6f284339f29c7030f7b092ebe5),
% lines 63233-65054.  Not a preserved v4.7 fragment; earlier layers are not rewritten.
% Content SHA256 of the carried block: 933b184650fc16b20e7f1330a8f0d63bdb02b372c49254be9001c1f06dad676b

% ============================================================================
\clearpage
\section*{v4.7.0-r85 Source-Only Citation Repair and Unity-Theorem Scope Separation}
\addcontentsline{toc}{section}{v4.7.0-r85 Source-Only Citation Repair and Unity-Theorem Scope Separation}
\label{sec:r85-citation-repair-unity-scope-separation}
% ============================================================================

\paragraph{Purpose.}
This r85 layer records two source repairs found by a systematic cross-reference audit
against the two cited source papers.  The first is a misdirected citation number; the second
is a scope separation inside the central unity theorem.  Neither repair adds a result, and
neither spends a guard.  Both narrow what is asserted; nothing previously proved is removed.

\noindent\textbf{Source-only status.}
This r85 layer is source-only.  It does not build a PDF, does not assign a new DOI, and does
not authorize a public release.  The intended reserved v4.7.0 Version DOI remains
\href{https://doi.org/10.5281/zenodo.20232911}{10.5281/zenodo.20232911}; the previous
baseline DOI \href{https://doi.org/10.5281/zenodo.20117361}{10.5281/zenodo.20117361}
remains provenance only.

\subsection*{r85 citation repair}
\label{subsec:r85-citation-repair}

Two passages cited the Stone--von Neumann step of \cite{brendecke2026rh} as Theorem~3.2.
In the cited source, the shared theorem counter makes Theorem~3.2 the functional equation
\(\xi(s)=\xi(1-s)\); the adelic Stone--von Neumann application is Step~3 in the proof of
Lemma~13.2, resting on the Weyl-algebra representation of Proposition~13.1.  The numbering is
confirmed internally: three further passages of this Companion cite Theorem~3.3 for the
Guinand--Weil explicit formula, which is the entry immediately following Theorem~3.2.

\[
\begin{array}{l|l|l}
\text{site} & \text{previous target} & \text{repaired target}\\
\hline
\text{proof of Theorem~\ref{thm:unity}} & \text{Theorem 3.2} &
\text{Lemma 13.2, Step 3}\\
\text{GUE universality argument} & \text{Theorem 3.2} &
\text{Proposition 13.1 and Lemma 13.2, Step 3}
\end{array}
\]

The second site requires both entries: Proposition~13.1 supplies the representation of the
adelic Weyl algebra \(\WA\) on \(\Hplusplus\), while Step~3 of Lemma~13.2 supplies
irreducibility and the resulting unitary equivalence.

\subsection*{r85 unity-theorem scope separation}
\label{subsec:r85-unity-scope-separation}

Theorem~\ref{thm:unity} previously listed the gauge structure
\(U(1)\times SU(2)\times SU(3)\) as item~(vi) of a list whose closing sentence asserts that
all items are unitarily equivalent by Stone--von Neumann.  That theorem relates irreducible
unitary representations of the Weyl algebra.  A Lie group is not such a representation, so the
equivalence sentence could not range over that item.

\paragraph{Theoremlet r85.}
Under the r60\(^+\) guarded grammar, the Stone--von Neumann equivalence of
Theorem~\ref{thm:unity} extends exactly over the items that are realized as irreducible
unitary representations of \(\WA\) on subspaces of \(\Hplus\):
\[
\mathsf{SvN}_{\rm scope}^{r85}=\{(\mathrm{i}),(\mathrm{ii}),(\mathrm{iii}),
(\mathrm{iv}),(\mathrm{v})\},
\qquad
U(1)\times SU(2)\times SU(3)\notin\mathsf{SvN}_{\rm scope}^{r85}.
\]
The gauge structure is retained in full as the symmetry group of those representations by
Theorem~\ref{thm:SM}, which is unchanged.

\paragraph{Proof sketch.}
Items (ii), (iii) and (v) are the three kernels and their sector reading, carried by
Theorem~\ref{thm:unitary-equiv}; item (iv) is the compensator image \(C(f)\in\Hplusplus\);
item (i) is the Sphaera, which \cite{brendecke2026rh} defines as a Hilbert space closed under
the generated dynamics, not as a bare geometric space.  Each is therefore a subspace of
\(\Hplus\) carrying the representation to which Lemma~13.2, Step~3 applies.  A Lie group
carries no such structure, so the equivalence cannot be asserted for it; the gauge statement
follows instead from the algebraic sector decomposition of Theorem~\ref{thm:SM}.  \(\square\)

\subsection*{r85 repair ledger}
\label{subsec:r85-repair-ledger}

The repairs close cleanly:
\[
R_{\rm rep}^{r85}=(0,0,0)
\]
for the three repaired items
\[
(\text{citation site 1},\;\text{citation site 2},\;\text{unity-theorem scope}).
\]
The repair is claim-reducing, not claim-adding.  No result is withdrawn: only the assertion
that the Stone--von Neumann equivalence ranges over the gauge group is withdrawn, while
Theorem~\ref{thm:SM} and every sector statement remain intact.

\paragraph{Divergence note.}
The same unity theorem also appears in the main paper \cite{brendecke2026main}, where it is
not modified by r85.  The two documents therefore differ in the scope sentence of that theorem
until the main paper is revised.  This is provenance information, not a claim change: the
present Companion carries the narrower and correct scope.

\subsection*{r85 guard update}
\label{subsec:r85-guard-update}

The r85 layer spends no guard and changes no residual:
\[
R_{\rm glob}^{(r85,+)}=(0,1,1,1,1),
\qquad
\Gamma_{\rm guard}^{r85}=(1,1,1,1,1),
\]
with the effective/local and exact/global ledgers unchanged from r84,
\[
R_{\rm eff/local}^{r85}=(0,0,0,0,0),
\qquad
R_{\rm exact/global}^{r85}=(1,1,1,1,1).
\]
The repairs do not imply exact finite \(C5\) zero closure, exact all-readout terminality,
\(C6/T2\), a global photon-clock/RH theorem, or global all-readout closure.

\subsection*{r85 release-path implication}
\label{subsec:r85-release-path-implication}

Because r85 corrects two citation targets and one scope sentence in the proof spine, any
later public-release package should be regenerated from r85 or later.  The r84 release-path
rule is otherwise unchanged: compact-titlepage discipline and final-handoff naming apply, and
no public release is authorized by r85 itself.


% ============================================================================
\clearpage
\section*{v4.7.0-r86 Source-Only: C5 Asymptotic-Collar Bookkeeping Consolidation}
\addcontentsline{toc}{section}{v4.7.0-r86 Source-Only: C5 Asymptotic-Collar Bookkeeping Consolidation}
\label{sec:r86-c5-asymptotic-collar-bookkeeping}
% ============================================================================

\paragraph{Purpose.}
This layer executes branch \(Q_5\) of the r48 queue.  The C5 collar facts are currently
spread over r04--r13 and the compressed r55--r80 chain, and the r83 table lists
\(C5_{\rm zero}\) as \(\mathsf{Open}_{\rm exact/global}\) with the reason
``\(\nexists N<\infty:\Phi_5^{(N)}=0\)''.  Consolidating the bookkeeping shows that this
entry mixes two different things: a question that the declared recursion answers, and a
question that it does not reach at all.  Separating them is the whole content of r86.  No
new mathematical result is produced; the gain is in classification.

\noindent\textbf{Source-only status.}
This r86 layer is source-only.  It does not build a PDF, does not assign a new DOI, and does
not authorize a public release.  The intended reserved v4.7.0 Version DOI remains
\href{https://doi.org/10.5281/zenodo.20232911}{10.5281/zenodo.20232911}.

\subsection*{r86 collar inventory}
\label{subsec:r86-collar-inventory}

The declared C5 collar data, collected:
\[
\begin{array}{l|l|l}
\text{object} & \text{content} & \text{origin}\\
\hline
\mathcal D_5^{\rm bd} & \text{bounded finite-stage collar domain} &
\text{Def.~\ref{def:v470-r09-bounded-c5-operator-domain}}\\
\mathcal U_5^{\rm bd} & \Phi\mapsto\tfrac34\Phi \text{ on }\mathcal D_5^{\rm bd} &
\text{Def.~\ref{def:v470-r09-bounded-c5-operator-domain}}\\
\text{admissibility} & \mathcal U_5^{\rm bd}\text{ preserves }\mathcal D_5^{\rm bd} &
\text{Lem.~\ref{lem:v470-r09-c5-finite-stage-admissibility}}\\
\Phi_5^{(n)} & (3/4)^n\text{ with }\Phi_5^{(0)}=1 &
\text{Thm.~\ref{thm:v470-r09-c5-bounded-operator-theorem}}\\
\text{finite positivity} & \Phi_5^{(n)}>0\ \forall n<\infty &
\text{Lem.~\ref{lem:v470-r11-finite-positive-guard}}\\
\mathsf{LG}_5^{\rm bd} & \text{limit may not replace a finite stage} &
\text{Def.~\ref{def:v470-r11-bounded-collar-limit-guard}}\\
\varepsilon\text{-use rule} & \varepsilon\text{-smallness allowed, promotion not} &
\text{Prop.~\ref{prop:v470-r11-asymptotic-use-guard}}\\
C5_{\rm det/eff} & \text{detector zero at }\theta=1/256,\ n=20 &
\text{r63--r64, compressed}
\end{array}
\]

Everything in this table is inherited.  r86 adds no entry to it.

\subsection*{r86 threshold function}
\label{subsec:r86-threshold-function}

The detector threshold is currently recorded only at the single point
\(\theta=1/256\), \(n=20\).  It is a value of a function, and writing that function down
costs nothing.

\paragraph{Theoremlet r86.1.}
For a detector threshold \(\theta\in(0,1)\), the least declared bounded stage at which the
collar score falls to or below \(\theta\) is
\[
n_{\min}(\theta)
=\left\lceil \log_{4/3}\!\frac1\theta \right\rceil
=\left\lceil \frac{\ln(1/\theta)}{\ln(4/3)} \right\rceil .
\]
In particular \(n_{\min}(1/256)=20\), and each halving of the threshold costs
\(\ln 2/\ln(4/3)=2.4094\ldots\) further stages.

\paragraph{Proof.}
By Theorem~\ref{thm:v470-r09-c5-bounded-operator-theorem},
\(\Phi_5^{(n)}=(3/4)^n\).  Since \(\ln(3/4)<0\), the condition
\((3/4)^n\le\theta\) is equivalent to \(n\ln(3/4)\le\ln\theta\), that is
\(n\ge\ln(1/\theta)/\ln(4/3)\); the least such integer is the stated ceiling.  For
\(\theta=2^{-8}\) this is \(\lceil 8\ln 2/\ln(4/3)\rceil=\lceil 19.276\ldots\rceil=20\),
which agrees with the inherited r63--r64 detector stage.  The halving cost is
\(\log_{4/3}2\). \(\square\)

\paragraph{Remark.}
The agreement at \(\theta=1/256\) is a consistency check on the inherited chain, not a new
claim: \((3/4)^{19}=0.004228\ldots>1/256\) while \((3/4)^{20}=0.003171\ldots\le 1/256\).

\subsection*{r86 finite-zero exclusion}
\label{subsec:r86-finite-zero-exclusion}

\paragraph{Theoremlet r86.2.}
Let \(\Phi^{(n+1)}=q\,\Phi^{(n)}\) be any purely multiplicative collar recursion with
\(q\in(0,1)\) and \(\Phi^{(0)}>0\).  Then
\[
\nexists\,N<\infty:\ \Phi^{(N)}=0 .
\]
The conclusion depends only on \(q>0\), not on the particular value \(q=3/4\).

\paragraph{Proof.}
\(\Phi^{(n)}=q^{n}\Phi^{(0)}\) is a product of finitely many strictly positive factors and is
therefore strictly positive for every finite \(n\). \(\square\)

\paragraph{Why this is worth recording.}
The proof is immediate and the statement is already implicit in
Remark~\ref{rem:v470-r09-not-c5-closure} and
Lemma~\ref{lem:v470-r11-finite-positive-guard}.  What has not been recorded is its
\emph{status consequence}.  ``No finite \(N\) with \(\Phi_5^{(N)}=0\)'' has been carried as
the reason for an \emph{open} entry, as though a finite closure might still be found by
working harder inside the collar.  Theoremlet r86.2 shows the opposite: inside the declared
recursion the answer is settled, and settled negatively.  Searching for a finite collar stage
that discharges C5 is not a hard open problem but a provably empty search.

The exclusion is also robust in a useful way.  Because it turns on \(q>0\) alone, no
refinement of the contraction factor can help: replacing \(3/4\) by any smaller positive
factor accelerates the decay and changes \(n_{\min}\), but leaves the exclusion untouched.

\subsection*{r86 status refinement}
\label{subsec:r86-status-refinement}

The r83 classes are \(\mathsf{Closed}_{\rm eff/local}\), \(\mathsf{Closed}_{\rm typed}\) and
\(\mathsf{Open}_{\rm exact/global}\).  The last of these conflates two situations.  We
separate them:
\[
\mathsf{Excluded}_{\rm rec}
\quad\text{versus}\quad
\mathsf{Open}_{\rm undecided}.
\]
\(\mathsf{Excluded}_{\rm rec}\) means: inside the declared structure the question has an
answer, and the answer is negative.  \(\mathsf{Open}_{\rm undecided}\) means: the declared
structure does not reach the question at all.

Applying this to \(C5_{\rm zero}\) splits it into two independent questions:
\[
\begin{array}{c|l|c}
& \text{question} & \text{r86 status}\\
\hline
(a) & \text{does }\mathcal U_5^{\rm bd}\text{ reach zero at a finite stage?} &
\mathsf{Excluded}_{\rm rec}\\
(b) & \text{is there a route to }C5=0\text{ outside the multiplicative collar?} &
\mathsf{Open}_{\rm undecided}
\end{array}
\]
Thus the r86 C5 decision datum is
\[
\mathsf D_{C5}^{r86}=\bigl(\mathsf{Excluded}_{\rm rec},\ \mathsf{Open}_{\rm undecided}\bigr).
\]

This is a refinement of the record, not a change of result.  \(C5=0\) remains unproved and
remains a forbidden assertion; question (b) is exactly the part that was never addressed, and
it is now stated as the open question in its own right rather than hidden inside (a).

\subsection*{r86 search direction}
\label{subsec:r86-search-direction}

Theoremlet r86.2 also says what a finite discharge would have to look like.  Since the
exclusion follows from multiplicativity together with positivity of the factor, any candidate
mechanism must leave that form.  The admissible shapes are:
\[
\begin{array}{c|l|l}
\text{shape} & \text{modification} & \text{what it would require}\\
\hline
\text{(i)} & \Phi^{(n+1)}=q\Phi^{(n)}-c,\ c>0 &
\text{a declared subtractive channel with its own guard}\\
\text{(ii)} & \Phi^{(n+1)}=0\text{ for }\Phi^{(n)}<\delta &
\text{a truncation rule; not admissible under }\mathsf{LG}_5^{\rm bd}\\
\text{(iii)} & \text{a non-collar object} &
\text{a separate witness, unrelated to }\mathcal U_5^{\rm bd}
\end{array}
\]
Shape (ii) is precisely what the r11 limit guard forbids, and r86 does not reopen it.
Shape (i) would need a subtractive term that is not currently declared anywhere in the
bounded--exterior--mixed grammar.  Shape (iii) is the honest reading of question (b): a
finite C5 discharge, if one exists, will not come from the collar iteration.

Nothing in r86 supplies any of the three.  The layer states the search space; it does not
enter it.

\subsection*{r86 ledger}
\label{subsec:r86-ledger}

The consolidation closes cleanly:
\[
R_{\rm cons}^{r86}=(0,0,0,0)
\]
for the four consolidated items
\[
(\text{inventory},\ \text{threshold function},\ \text{finite-zero exclusion},
\ \text{status split}).
\]
The exact/global ledger is unchanged, because question (b) remains unanswered:
\[
R_{\rm exact/global}^{r86}=(1,1,1,1,1).
\]
The first component keeps the value \(1\), now with a sharper reason: not ``no finite
closure has been found'' but ``no finite closure exists inside the declared recursion, and
no route outside it has been declared''.

\subsection*{r86 guard update}
\label{subsec:r86-guard-update}

r86 spends no guard.  It strengthens the justification of the standing C5 guard rather than
consuming it:
\[
R_{\rm glob}^{(r86,+)}=(0,1,1,1,1),
\qquad
\Gamma_{\rm guard}^{r86}=(1,1,1,1,1),
\qquad
R_{\rm eff/local}^{r86}=(0,0,0,0,0).
\]
The layer does not imply exact finite \(C5\) zero closure, exact all-readout terminality,
\(C6/T2\), a global photon-clock/RH theorem, or global all-readout closure.  The
detector-threshold result \(C5_{\rm det/eff}\) is unaffected and remains
\(\mathsf{Closed}_{\rm eff/local}\) at \(\theta=1/256\), \(n=20\).

\subsection*{r86 queue update}
\label{subsec:r86-queue-update}

Branch \(Q_5\) of the r48 queue is discharged as a bookkeeping task.  The remaining
guard-neutral branch is
\[
Q_3:\ \text{photon-clock selector-window calibration},
\]
which would sharpen the provenance of the bound \(\Delta_\gamma<1/5\) carried by
\(W_{\rm clk}^{\rm loc}\).  \(Q_4\) remains inadmissible while \(\beta_4=1/256\) is to stay
unspent.


% ============================================================================
\clearpage
\section*{v4.7.0-r87 Source-Only: Photon-Clock Selector-Window Calibration}
\addcontentsline{toc}{section}{v4.7.0-r87 Source-Only: Photon-Clock Selector-Window Calibration}
\label{sec:r87-photon-clock-window-calibration}
% ============================================================================

\paragraph{Purpose.}
This layer executes branch \(Q_3\) of the r48 queue.  The r15 finite window family is
recorded as a table of three computed values; the rules behind that table are not written
down in closed form.  r87 extracts them, verifies the inherited numbers, and reduces the one
remaining blocked relay edge to a single missing datum.  As in r86, no new mathematical
result is produced: everything below follows from the r15 extraction rule and the r16 lift
computation.

\noindent\textbf{Source-only status.}
This r87 layer is source-only.  It does not build a PDF, does not assign a new DOI, and does
not authorize a public release.  It makes no spectral identification between Riemann zeros
and photon ticks and does not touch the \(W_\gamma^{\rm glob}\) guard.

\subsection*{r87 window inventory}
\label{subsec:r87-window-inventory}

The inherited photon-clock window data, from
Section~\ref{sec:v470-r15-zero-clock-phase-window-relay-stitching}:
\[
\begin{array}{l|l}
\text{overlap flux} & S_V=N_V+\kappa_V N_V(N_V-1)\\
\text{required gain} & G_V^{\rm req}=16/S_V,\quad 16=\Theta_v/a^2
 \text{ at }\Theta_v=1/256,\ a=1/64\\
\text{relay criterion} & G_VS_V\ge16\\
\text{window} & I_V=\{\tau:\mathcal C_\gamma(\tau)\ge G_V^{\rm req}\}\\
\text{half-width} & \delta_V=\sigma_V\arccos\sqrt{(G_V^{\rm req}-1)/(G_V^{\rm avail}-1)}
\end{array}
\]

\subsection*{r87 window-existence criterion}
\label{subsec:r87-window-existence}

\paragraph{Theoremlet r87.1.}
A chart \(V\) has a nonempty zero-clock phase window exactly when its overlap flux lies in
the half-open interval
\[
\frac{16}{G_V^{\rm avail}}\ \le\ S_V\ <\ 16 .
\]
The upper end is the direct/assisted boundary: for \(S_V\ge16\) the chart needs no clock
window at all.  The lower end is the assisted/blocked boundary.

\paragraph{Proof.}
The r15 extraction condition is \(1<G_V^{\rm req}\le G_V^{\rm avail}\) with
\(G_V^{\rm req}=16/S_V\).  The left inequality \(1<16/S_V\) is equivalent to \(S_V<16\); the
right inequality \(16/S_V\le G_V^{\rm avail}\) is equivalent to
\(S_V\ge16/G_V^{\rm avail}\). \(\square\)

\paragraph{Verification against the inherited table.}
\[
\begin{array}{c|c|c|c|c}
\text{chart} & S_V & \text{admissible band} & \text{in band} & \delta_V/\sigma_V\\
\hline
V3 & 10 & [8,16) & \text{yes} & 0.684719\\
V5 & 10 & [32/5,16) & \text{yes} & 0.886077\\
V8 & 9 & [8,16) & \text{yes} & 0.490883
\end{array}
\]
The three half-widths were recomputed from the r15 formula in exact rational ratios
\(3/5\), \(2/5\), \(7/9\) and agree with the inherited values to six decimals.

\paragraph{Remark.}
The half-width is monotone in both arguments: it tends to \(\pi/2\,\sigma_V\) as
\(G_V^{\rm req}\to1^+\), and it closes to zero as \(G_V^{\rm req}\to G_V^{\rm avail}\).  A
chart sitting at the lower band edge therefore has a window in principle but no timing margin
in practice.  Of the three extracted charts, \(V8\) is closest to that edge.

\subsection*{r87 single-packet lift formula}
\label{subsec:r87-single-packet-lift}

The r16 shadow-basin lift moved \(N_6:3\to4\) at fixed \(\kappa_6=3/4\), \(G_6=3/2\), which
raised \(S_6\) from \(15/2\) to \(13\).  That step is an instance of a general rule.

\paragraph{Theoremlet r87.2.}
Adding one admissible packet at fixed \(\kappa\) changes the overlap flux by
\[
\Delta S(N\to N+1)=1+2\kappa N .
\]

\paragraph{Proof.}
\(S(N+1)-S(N)=\bigl[(N+1)+\kappa(N+1)N\bigr]-\bigl[N+\kappa N(N-1)\bigr]=1+2\kappa N\).
\(\square\)

For the r16 lift this gives \(1+2\cdot\tfrac34\cdot3=\tfrac{11}{2}\), and indeed
\(15/2+11/2=13\).  The gain is quadratic in \(N\): a single packet is worth more at a
populated front than at a sparse one.

\subsection*{r87 blocked-chart ranking}
\label{subsec:r87-blocked-ranking}

r15 ranks the blocked charts by the gain ratio \(G^{\rm req}/G^{\rm avail}\) and concludes
that \(V6\) is closest.  That ranking is correct for its own measure, but it is not the only
measure, and the two do not agree:
\[
\begin{array}{c|c|c|c|c}
\text{chart} & S_V & G_V^{\rm avail} &
G^{\rm req}/G^{\rm avail} & \text{missing }S\\
\hline
V0 & 4 & 1 & 4 & 12\\
V6\ (\text{pre-r16}) & 15/2 & 3/2 & 64/45\approx1.422 & 19/6\approx3.167\\
V9\ (\text{pre-r17}) & 3 & 3 & 16/9\approx1.778 & 7/3\approx2.333
\end{array}
\]
By gain ratio the order is \(V6<V9<V0\); by missing overlap flux it is \(V9<V6<V0\).  The
two measures correspond to two different lift mechanisms: raising the available clock gain
acts multiplicatively, adding packets acts through \(S\).  Both r16 and r17 in fact proceeded
through \(S\), by adding a single packet, so the additive column is the one that governed the
executed moves.  It also explains why \(V9\) yielded to the smaller lift: it needed \(7/3\)
where \(V6\) needed \(19/6\).

\subsection*{r87 lift formula against the recorded V9 micro-lift}
\label{subsec:r87-v9-reduction}

The chart parameters are recorded in the r14 scan ledger, which lists
\((N_9,\kappa_9,G_9)=(2,\tfrac12,3)\), hence \(S_9=3\) and \(G_9S_9=9<16\).  The blocked
edge \(V8\to V9\) needs \(\Delta S\ge7/3\) by the ranking table above, and Theoremlet r87.2
gives for a single added packet
\[
\Delta S=1+2\kappa_9N_9=1+2\cdot\tfrac12\cdot2=3\ \ge\ \tfrac73 .
\]
So one packet suffices, with \(S_9'=6\) and \(G_9S_9'=18\ge16\).

This is exactly the move carried out in
Section~\ref{sec:v470-r17-atlas-gap-bridge-v9-microlift}: the r17 guarded micro-lift takes
\((2,\tfrac12,3)\mapsto(3,\tfrac12,3)\) and records \(H_9'=18>16\) with margin
\(\mu_9'=2\).  The r87 formula therefore reproduces an already executed lift from the chart
parameters alone, without reference to the r17 computation.

\paragraph{Status of the chain.}
After the r16 lift of \(V6\), the r17 atlas-gap bridge \(V6\to V7\) and the r17 micro-lift of
\(V9\), the finite chart chain runs \(V1\to\cdots\to V9\).  The only chart outside it is
\(V0\), which needs \(\Delta S=12\) --- four times the largest lift so far executed --- and
which r15--r17 keep explicitly guarded.  The calibration confirms that ordering: by the
additive measure \(V0\) is far outside the range in which single-packet moves operate.

\subsection*{r87 provenance gap for \texorpdfstring{\(\Delta_\gamma<1/5\)}{Delta gamma < 1/5}}
\label{subsec:r87-provenance-gap}

The stated aim of \(Q_3\) was to sharpen the provenance of the bound \(\Delta_\gamma<1/5\)
carried by \(W_{\rm clk}^{\rm loc}\).  That bound occurs in the present source only inside
summary and guard layers, never with a derivation: it enters through the compressed
r55--r80 chain, where the detailed r65 step has been replaced by the chain invariant.  The
uncompressed log is not part of this source packet.

Accordingly r87 does \emph{not} discharge that part of \(Q_3\).  What can be said from the
present source is negative and worth recording:
\[
\Delta_\gamma<1/5 \notin \{\text{quantities derivable from the r15--r16 window layer}\}.
\]
The r15 layer produces the ratios \(3/5\), \(2/5\), \(7/9\), the gains \(8/5\), \(16/9\),
\(32/15\), \(16/3\) and the threshold constant \(16\); none of these yields \(1/5\), and the
window half-widths are measured in units of \(\sigma_V\) rather than as a bound on a
frequency spread.  The bound therefore belongs to a different layer and must be recovered
from the uncompressed r65 text, not reconstructed from the finite chart family.

\subsection*{r87 ledger}
\label{subsec:r87-ledger}

The calibration closes in three of four items:
\[
R_{\rm cal}^{r87}=(0,0,0,1)
\]
for
\[
(\text{window-existence criterion},\ \text{lift formula},\ \text{blocked-chart ranking},
\ \Delta_\gamma\text{ provenance}).
\]
The fourth component stays at \(1\): the bound \(\Delta_\gamma<1/5\) remains an inherited
constant whose derivation is outside this source.  This is a recorded documentation gap, not
a mathematical obstruction, and it does not affect the standing of
\(W_{\rm clk}^{\rm loc}\), which is carried as a conditional local statement.

\subsection*{r87 guard update}
\label{subsec:r87-guard-update}

r87 spends no guard:
\[
R_{\rm glob}^{(r87,+)}=(0,1,1,1,1),
\qquad
\Gamma_{\rm guard}^{r87}=(1,1,1,1,1),
\qquad
R_{\rm eff/local}^{r87}=(0,0,0,0,0).
\]
All statements are finite, local and chart-bound.  The layer does not identify Riemann zeros
with photon ticks, does not prove a global photon-clock/RH theorem, does not prove
\(C6/T2\), and does not remove the all-readout blocker.  The relay chains remain finite:
the longest certified chain is the length-six chain of
Corollary~\ref{cor:v470-r16-finite-chain-status}.

\subsection*{r87 queue update}
\label{subsec:r87-queue-update}

Branch \(Q_3\) is discharged in its calibration part and left open in its provenance part.
The r48 queue after r87:
\[
\begin{array}{c|l|l}
\text{branch} & \text{content} & \text{status}\\
\hline
Q_1 & \text{local field-interface elaboration} & \text{done (r49--r54)}\\
Q_2 & \text{theoremlet-to-lemma packaging} & \text{done (r55--r80)}\\
Q_3 & \text{photon-clock selector-window calibration} &
\text{calibrated (r87); }\Delta_\gamma\text{ provenance open}\\
Q_4 & \text{C4 mixed-kernel localization} & \text{inadmissible, }\beta_4\text{ unspent}\\
Q_5 & \text{C5 asymptotic-collar bookkeeping} & \text{done (r86)}
\end{array}
\]
One concrete item remains that needs no new guard: the uncompressed r65 text carrying the
derivation of \(\Delta_\gamma<1/5\).


% ============================================================================
\clearpage
\section*{v4.7.0-r88 Source-Only: Admissibility Test for Non-Collar Routes to Finite C5 Discharge}
\addcontentsline{toc}{section}{v4.7.0-r88 Source-Only: Admissibility Test for Non-Collar Routes to Finite C5 Discharge}
\label{sec:r88-non-collar-route-admissibility}
% ============================================================================

\paragraph{Purpose.}
r86 split the \(C5_{\rm zero}\) entry into question (a), settled negatively inside the
declared multiplicative collar recursion, and question (b): whether a route to a finite
\(C5\) discharge exists \emph{outside} that recursion.  r86 listed three candidate shapes
without examining them.  This layer examines them.  Two are eliminated; the third is not,
and question (b) remains open.  Nothing here asserts that a finite discharge exists.

\noindent\textbf{Source-only status.}
This r88 layer is source-only.  It does not build a PDF, does not assign a new DOI, and does
not authorize a public release.  The forbidden assertion \(C5=0\) is not made, and no guard
is spent.  The analysis is conditional throughout: it asks what a candidate mechanism would
have to satisfy, not what holds.

\subsection*{r88 subtractive channel: closed form}
\label{subsec:r88-subtractive-closed-form}

Shape (i) of the r86 table is the constant-subtraction recursion
\[
\Phi^{(n+1)}=q\,\Phi^{(n)}-c,
\qquad q\in(0,1),\ c>0,\ \Phi^{(0)}>0 .
\]

\paragraph{Theoremlet r88.1.}
This recursion has the closed form
\[
\Phi^{(n)}
=q^{n}\!\left(\Phi^{(0)}+\frac{c}{1-q}\right)-\frac{c}{1-q},
\]
and it attains the value zero at an integer stage \(N\) if and only if
\[
c=c_\star(N):=\frac{(1-q)\,q^{N}\,\Phi^{(0)}}{1-q^{N}} .
\]
For \(q=3/4\), \(\Phi^{(0)}=1\) this is
\[
c_\star(N)=\frac{1}{4\bigl((4/3)^{N}-1\bigr)} .
\]

\paragraph{Proof.}
Writing \(\psi^{(n)}=\Phi^{(n)}+c/(1-q)\) turns the recursion into
\(\psi^{(n+1)}=q\psi^{(n)}\), whence \(\psi^{(n)}=q^n\psi^{(0)}\) and the displayed form.
Setting \(\Phi^{(N)}=0\) gives \(q^{N}(\Phi^{(0)}+c/(1-q))=c/(1-q)\); solving for \(c\)
yields \(c_\star(N)\).  Conversely \(c=c_\star(N)\) returns \(\Phi^{(N)}=0\) by substitution.
\(\square\)

\paragraph{Verified instances.}
For \(q=3/4\), \(\Phi^{(0)}=1\): \(c_\star(1)=3/4\), \(c_\star(2)=9/28\),
\(c_\star(3)=27/148\), \(c_\star(4)=81/700\), \(c_\star(5)=243/3124\).  Each was checked by
running the recursion in exact rational arithmetic; each reaches zero at its own stage and at
no earlier one.

\paragraph{Remark on genericity.}
For \(c\ne c_\star(N)\) for every \(N\), the orbit does not meet zero at all: it passes from
positive to negative between two consecutive stages.  Sample runs with
\(c=1/10,1/5,1/4,3/10\) cross zero at stages \(5,3,3,3\) respectively without ever taking the
value zero.  Exact finite discharge by a constant subtractive channel therefore requires the
channel strength to be tuned to the target stage --- a measure-zero condition on \(c\), not a
structural feature.

\subsection*{r88 the admissibility obstruction}
\label{subsec:r88-admissibility-obstruction}

Tuning is not the decisive objection.  The decisive objection is what happens immediately
afterwards.

\paragraph{Theoremlet r88.2.}
Let \(c=c_\star(N)\), so that \(\Phi^{(N)}=0\).  Then
\[
\Phi^{(N+1)}=-c<0 ,
\]
and \(\Phi^{(n)}<0\) for all \(n>N\).  Hence the orbit leaves the bounded collar domain
\(\mathcal D_5^{\rm bd}\) at the stage immediately after the discharge.

\paragraph{Proof.}
\(\Phi^{(N+1)}=q\cdot0-c=-c<0\).  For \(n>N\) the closed form of Theoremlet r88.1 is
strictly decreasing once negative, since \(q\Phi-c<\Phi\) whenever \(\Phi<c/(1-q)\), which
holds for all \(\Phi\le0\).  Admissibility of \(\mathcal U_5^{\rm bd}\) on
\(\mathcal D_5^{\rm bd}\) was established in
Lemma~\ref{lem:v470-r09-c5-finite-stage-admissibility} precisely through preservation of
positivity; a map whose next iterate is negative fails that condition. \(\square\)

\paragraph{Corollary r88.3 (shape (i) collapses into shape (ii)).}
A subtractive channel can deliver a finite discharge only if it is accompanied by a stopping
rule that halts the recursion at stage \(N\) instead of continuing into the negative range.
Such a rule is a truncation, which is shape (ii) of the r86 table, and truncation is
forbidden by the bounded-collar limit guard
\(\mathsf{LG}_5^{\rm bd}\) of
Definition~\ref{def:v470-r11-bounded-collar-limit-guard}.  Shapes (i) and (ii) are therefore
not two independent candidates but one, and it is already excluded by a standing guard.

\subsection*{r88 the free-parameter trap}
\label{subsec:r88-free-parameter-trap}

One might weaken shape (i) to a stage-dependent channel
\(\Phi^{(n+1)}=q\Phi^{(n)}-c_n\).  This does produce a finite zero: choosing
\(c_N=q\Phi^{(N)}\) sets \(\Phi^{(N+1)}=0\) for any prescribed \(N\).

\paragraph{Remark r88.4.}
That freedom is the objection, not the solution.  A family with an unconstrained sequence
\((c_n)\) can realize \emph{any} prescribed orbit whatsoever, including orbits that violate
every other guard in the ledger.  A mechanism that can produce the desired conclusion and its
negation with equal ease carries no information about \(C5\); it records the desired answer
rather than deriving it.  An admissible stage-dependent channel would therefore have to fix
\((c_n)\) by an independent structural rule --- and it is that rule, not the resulting zero,
that would constitute the mathematical content.  No such rule is declared anywhere in the
bounded--exterior--mixed grammar.

\subsection*{r88 what remains: shape (iii)}
\label{subsec:r88-shape-three}

Shape (iii) is the only surviving candidate: an object unrelated to the collar iteration.
r88 does not supply one.  What it can do is state the burden such an object would carry.

\[
\begin{array}{c|l}
& \text{requirement on a shape-(iii) witness}\\
\hline
(\alpha) & \text{it must exhibit the vanishing of the }C5\text{ obstruction itself,}\\
 & \text{not the smallness of the score }\Phi_5\\
(\beta) & \text{it must be independent of }\mathcal U_5^{\rm bd}\text{, since by r86 that
 recursion}\\
 & \text{cannot deliver the conclusion}\\
(\gamma) & \text{it must remain compatible with }C5_{\rm det/eff}\text{, which is already}\\
 & \text{established at }\theta=1/256,\ n=20\\
(\delta) & \text{it must not import the all-readout slot, since }0_{\rm termdata}
 \not\Rightarrow0_{\rm all}\\
 & \text{remains frozen}
\end{array}
\]

Requirement \((\alpha)\) is the substantive one.  The collar layer measures \(C5\) through a
score whose finite values are positive and whose limit is zero.  Any statement about \(C5\)
obtained \emph{through that score} is therefore a statement about the limit, which is exactly
what the standing guard separates from finite discharge.  A shape-(iii) route must reach the
obstruction by some other access; otherwise it re-enters the collar layer and inherits its
verdict.

\subsection*{r88 ledger}
\label{subsec:r88-ledger}

Of the three r86 shapes, two are now closed and one is not:
\[
\mathsf D_{(b)}^{r88}
=\bigl(\text{(i) excluded},\ \text{(ii) excluded},\ \text{(iii) open}\bigr),
\qquad
R_{(b)}^{r88}=(0,0,1).
\]
Question (b) of the r86 split therefore remains \(\mathsf{Open}_{\rm undecided}\), but its
search space is narrower: it is no longer ``some mechanism outside the collar'' but
specifically a witness meeting \((\alpha)\)--\((\delta)\), with the two parametric
modifications of the collar recursion ruled out.

\subsection*{r88 guard update}
\label{subsec:r88-guard-update}

r88 spends no guard.  It closes two candidate routes and leaves every ledger value where it
stood:
\[
R_{\rm glob}^{(r88,+)}=(0,1,1,1,1),
\qquad
\Gamma_{\rm guard}^{r88}=(1,1,1,1,1),
\qquad
R_{\rm exact/global}^{r88}=(1,1,1,1,1).
\]
Corollary r88.3 strengthens \(\mathsf{LG}_5^{\rm bd}\) by showing that the most natural way
around it collapses back into it.  The layer does not imply exact
finite \(C5\) zero closure, exact all-readout terminality, \(C6/T2\), a global
photon-clock/RH theorem, or global all-readout closure.


% ============================================================================
\clearpage
\section*{v4.7.0-r89 Source-Only: r87 Correction --- Recorded V9 Data and Executed Micro-Lift}
\addcontentsline{toc}{section}{v4.7.0-r89 Source-Only: r87 Correction --- Recorded V9 Data and Executed Micro-Lift}
\label{sec:r89-r87-correction}
% ============================================================================

\paragraph{Purpose.}
The r87 layer, as first written, asserted that the decomposition of the \(V9\) overlap flux
was not recorded in this source and that \(V8\to V9\) was the last blocked relay edge.  Both
assertions were wrong.  This layer records the correction and its cause.

\noindent\textbf{Source-only status.}
This r89 layer is source-only.  It does not build a PDF, does not assign a new DOI, and does
not authorize a public release.

\subsection*{r89 what was wrong}
\label{subsec:r89-what-was-wrong}

\[
\begin{array}{l|l|l}
\text{r87 assertion} & \text{status} & \text{actual location}\\
\hline
\text{``}(N_9,\kappa_9)\text{ not recorded''} & \text{false} &
\text{r14 chart ledger: }(2,\tfrac12,3)\\
\text{``}V8\to V9\text{ is the only blocked edge''} & \text{false} &
\text{r17 stitched it}\\
\text{``no claim which case holds''} & \text{moot} &
\text{decided in r17}
\end{array}
\]

The r14 finite valley-front chart scan lists all ten charts with packet count, coherence
parameter and available gain; \(V9\) appears there as \((N,\kappa,G)=(2,\tfrac12,3)\) with
clock score \(9/16\) and status blocked.  The r17 guarded micro-lift then performed
\((2,\tfrac12,3)\mapsto(3,\tfrac12,3)\), obtaining \(H_9'=18>16\) with margin
\(\mu_9'=2\), and stitched the chain to \(V1\to\cdots\to V9\).  Both facts stand in the same
document as r87.

\subsection*{r89 what was right}
\label{subsec:r89-what-was-right}

The defect was in the source survey, not in the calculation.  The r87 single-packet lift
formula
\[
\Delta S=1+2\kappa N
\]
applied to the r14 parameters gives \(\Delta S=3\), hence \(S_9'=6\) and \(G_9S_9'=18\) ---
precisely the values r17 reports.  The formula therefore reproduces the executed micro-lift
from the chart parameters alone.  What r87 presented as an open alternative was in fact a
correct prediction of a step already taken.

The corrected r87 text records this as a verification rather than as a question.  The r87
ledger is unchanged: \(R_{\rm cal}^{r87}=(0,0,0,1)\), whose single open component was always
the \(\Delta_\gamma\) provenance and never the \(V9\) datum.

\subsection*{r89 cause and rule}
\label{subsec:r89-cause-and-rule}

The survey behind r87 read r15 and r16 and inferred from the absence of the decomposition
\emph{there} that it was absent from the document.  It did not read r14, one round earlier,
where the chart ledger is tabulated, nor r17, one round later, where the lift is executed.

\paragraph{Rule r89.}
A claim that a datum is missing from this source is a claim about the whole monolith and must
be checked against the rounds adjacent to the one being read, not only against the round in
which the datum is used.  Absence in the round of use is evidence of nothing: the solve chain
habitually tabulates in one round and consumes in another.

\subsection*{r89 ledger and provenance}
\label{subsec:r89-ledger-provenance}

The correction is editorial and touches no result:
\[
R_{\rm corr}^{r89}=(0,0,0)
\]
for (chart datum restored, chain status restored, prediction reclassified as verification).

RC22 was packaged before the correction and therefore carries the erroneous r87 text.  It is
retained unchanged as provenance; the present source and any later release candidate carry
the corrected text.  No earlier round is modified by r89.

\subsection*{r89 guard update}
\label{subsec:r89-guard-update}

r89 spends no guard and changes no ledger value:
\[
R_{\rm glob}^{(r89,+)}=(0,1,1,1,1),
\qquad
\Gamma_{\rm guard}^{r89}=(1,1,1,1,1).
\]
The finite relay chain \(V1\to\cdots\to V9\) is inherited from r17 and is not extended here.
\(V0\) remains outside it and guarded.  The layer does not imply exact finite \(C5\) zero
closure, exact all-readout terminality, \(C6/T2\), a global photon-clock/RH theorem, or
global all-readout closure.


% ============================================================================
\clearpage
\section*{v4.7.0-r90 Source-Only: Reconstruction Attempt for the Bound \texorpdfstring{\(\Delta_\gamma<1/5\)}{Delta gamma < 1/5}}
\addcontentsline{toc}{section}{v4.7.0-r90 Source-Only: Reconstruction Attempt for the Bound Delta gamma < 1/5}
\label{sec:r90-delta-gamma-reconstruction}
% ============================================================================

\paragraph{Purpose and status.}
r87 recorded that the bound \(\Delta_\gamma<1/5\) carried by \(W_{\rm clk}^{\rm loc}\) has no
derivation in this source: it enters through the compressed r55--r80 chain, and the detailed
r65 text is not part of the source packet.  This layer attempts a reconstruction from the
material that \emph{is} present.  It arrives at a bound of the right size from the r15 clock
model.  \textbf{This is a reconstruction, not a recovery.}  It does not establish that r65
argued this way, and it does not replace the missing text.  The provenance component of
\(R_{\rm cal}^{r87}\) stays open.

\subsection*{r90 what the source fixes}
\label{subsec:r90-what-is-fixed}

Three ingredients are stated in
Section~\ref{sec:v470-r15-zero-clock-phase-window-relay-stitching}.  The photon-clock probe is
\[
\mathcal C_\gamma(\tau)
=1+\eta\,
\frac{\bigl|\sum_{m=1}^{M}w_me^{i\gamma_m\tau}\bigr|^2}
{\bigl(\sum_{m=1}^{M}|w_m|\bigr)^2},
\]
its local normalized form is
\[
G_V(\tau)=1+(G_V^{\rm avail}-1)\cos^2\!\Bigl(\frac{\tau-\tau_V}{\sigma_V}\Bigr),
\]
so that \(\eta=G_V^{\rm avail}-1\), and window membership requires
\(\mathcal C_\gamma(\tau)\ge G_V^{\rm req}\).

\subsection*{r90 coherence bound}
\label{subsec:r90-coherence-bound}

\paragraph{Theoremlet r90.}
Let the clock phases be \(\varphi_m=\gamma_m\tau\), centred so that
\(\langle\varphi\rangle=0\), and write \(\mathrm{Var}(\varphi)\) for their weighted variance.
To first order in the spread,
\[
\frac{\bigl|\sum_m w_me^{i\varphi_m}\bigr|^2}{\bigl(\sum_m|w_m|\bigr)^2}
=1-\mathrm{Var}(\varphi)+O\bigl(\mathrm{Var}^2\bigr),
\]
and the selector is preserved on chart \(V\) exactly while
\[
\mathrm{Var}(\varphi)\ \le\ 1-\frac{G_V^{\rm req}-1}{G_V^{\rm avail}-1}
\ =\ 1-\rho_V ,
\]
where \(\rho_V\) is the ratio already tabulated in the r15 half-width formula.

\paragraph{Proof.}
Expanding \(\langle e^{i\varphi}\rangle=1+i\langle\varphi\rangle-\tfrac12\langle\varphi^2\rangle
+O(\varphi^3)\) with \(\langle\varphi\rangle=0\) gives
\(\langle e^{i\varphi}\rangle=1-\tfrac12\mathrm{Var}+O(\mathrm{Var}^2)\), hence
\(|\langle e^{i\varphi}\rangle|^2=1-\mathrm{Var}+O(\mathrm{Var}^2)\).  Substituting into
\(\mathcal C_\gamma\ge G^{\rm req}\) with \(\eta=G^{\rm avail}-1\) yields
\(1+(G^{\rm avail}-1)(1-\mathrm{Var})\ge G^{\rm req}\), which rearranges to the display.
\(\square\)

\subsection*{r90 the resulting numerical bound}
\label{subsec:r90-numerical-bound}

Applied to the three charts that actually carry windows:
\[
\begin{array}{c|c|c|c}
\text{chart} & \rho_V & 1-\rho_V & \text{decimal}\\
\hline
V3 & 3/5 & 2/5 & 0.4000\\
V5 & 2/5 & 3/5 & 0.6000\\
V8 & 7/9 & 2/9 & 0.2222
\end{array}
\]
A bound valid for the whole family is the minimum,
\[
\Delta_\gamma^{\max}=\min_V\bigl(1-\rho_V\bigr)=\frac29=0.2222\ldots,
\]
attained at the narrowest window \(V8\).  The recorded bound \(1/5=0.2\) sits just below it,
with a margin of ten percent:
\[
\frac15<\frac29,
\qquad
1-\frac{1/5}{2/9}=\frac1{10}.
\]

\paragraph{Assessment.}
A conservative rounding of \(2/9\) to the nearest simple fraction below is \(1/5\), and that
matches the rounding discipline used elsewhere in this document, where thresholds are carried
as \(1/256\) and \(1/64\) rather than as decimals.  The reconstruction therefore produces a
bound of the correct magnitude, at the correct place in the argument, with a plausible safety
margin.

\subsection*{r90 rejected alternative}
\label{subsec:r90-rejected-alternative}

One numerical coincidence in the same table should be recorded and dismissed:
\[
\rho_{V3}-\rho_{V5}=\frac35-\frac25=\frac15 .
\]
This reproduces the constant exactly rather than approximately, but it carries no mechanism:
a difference of two window ratios is not a spread of clock frequencies, it does not depend on
\(\gamma\) at all, and nothing in the grammar would make such a difference a preservation
threshold.  It is rejected as coincidence.

\subsection*{r90 limits of the reconstruction}
\label{subsec:r90-limits}

Three gaps remain and are stated rather than papered over:
\[
\begin{array}{c|l}
(\alpha) & \Delta_\gamma\text{ is nowhere defined in this source; whether it denotes a
 variance,}\\
 & \text{a standard deviation or a maximal deviation is not recorded}\\
(\beta) & \text{the expansion is first order and holds only for small spread}\\
(\gamma) & \text{the ten-percent margin is plausible but not derived}
\end{array}
\]
Under reading \((\alpha)\) as a standard deviation the bound would instead compare against
\(\sqrt{2/9}=0.471\), and \(1/5\) would be far more conservative than necessary; under the
variance reading it is a close fit.  That is evidence for the variance reading, not proof.

\subsection*{r90 ledger and guard update}
\label{subsec:r90-ledger-guard}

The reconstruction is recorded as such:
\[
R_{\rm rec}^{r90}=(0,0,1)
\]
for (coherence bound derived, numerical match established, r65 provenance still missing).
The third component stays at \(1\): only the original text can close it.

r90 spends no guard:
\[
R_{\rm glob}^{(r90,+)}=(0,1,1,1,1),
\qquad
\Gamma_{\rm guard}^{r90}=(1,1,1,1,1).
\]
All statements are local and chart-bound.  The layer makes no spectral identification between
Riemann zeros and photon ticks, and it does not imply exact finite \(C5\) zero closure, exact
all-readout terminality, \(C6/T2\), a global photon-clock/RH theorem, or global all-readout
closure.  \(W_{\rm clk}^{\rm loc}\) continues to be carried as a conditional local statement
under its inherited bound, whose status as an inherited constant is unchanged by r90.


% ============================================================================
\clearpage
\section*{v4.7.0-r91 Source-Only: Q4 Executed --- Why \texorpdfstring{\(\beta_4\)}{beta 4} Cannot Be Spent}
\addcontentsline{toc}{section}{v4.7.0-r91 Source-Only: Q4 Executed --- Why beta 4 Cannot Be Spent}
\label{sec:r91-q4-beta4-not-spendable}
% ============================================================================

\paragraph{Purpose.}
Branch \(Q_4\) of the r48 queue has been carried as \emph{inadmissible} since r54, with the
reason recorded as ``\(\beta_4=1/256\) remains unspent''.  That reads as a policy: a guard
held back by choice.  This layer executes the branch and finds that it is not a policy.  The
attempt terminates for a structural reason already present in the source, and the guard
survives not because it is protected but because spending it would be the very operation the
matter-frame typing forbids.  \textbf{No guard is spent here.}

\noindent\textbf{Source-only status.}
This r91 layer is source-only.  It does not build a PDF, does not assign a new DOI, and does
not authorize a public release.

\subsection*{r91 what C4 still is}
\label{subsec:r91-what-c4-is}

The first finding is that \(C4\) is no longer a commutator problem.  The r07 split test
(Lemma~\ref{lem:v470-r07-commutator-component-routed}) decomposes
\[
S_4^{\rm comm}\leadsto S_{4,-}^{\rm comm}\oplus S_{4,+}^{\rm comm},
\qquad
S_{4,-}^{\rm comm}:\mathsf{routed}_{\rm comm},
\quad
S_{4,+}^{\rm comm}:\mathsf{zero}_{2s},
\]
and neither component carries active leakage afterwards.  The antisymmetric mixed-transport
part is not contracted away; it is \emph{read through a declared commutator comparison
channel}, and the source is explicit that this ``does not mean a global commutator vanishing
theorem''.  The active support reduces to a single class,
\[
\operatorname{Supp}_4^{r07}=\{S_{4,\rm out}^{\rm collar}\},
\]
the component that crosses the bounded--exterior interface.  What remains of \(C4\) is
therefore a boundary-crossing seam, not a commutator remainder.

\subsection*{r91 what r13 already settled}
\label{subsec:r91-what-r13-settled}

Lemma~\ref{lem:v470-r13-nonabsorption-test} tests absorbability of that residual kernel and
finds it negative: it is not absorbable by the declared bounded, comparison-routed or
selector-routed channels.  The same layer enumerates what an absorption would require --- a
new declared comparison tensor, a new selector, or an identification with bounded \(C5\)
contraction debt --- and records that the first two are absent from the datum while the third
is forbidden by the separation rule.

So \(Q_4\) cannot proceed by absorption inside the present grammar.  That much was known.
The question r91 asks is different: what would it \emph{mean} to spend \(\beta_4\)?

\subsection*{r91 the threshold identity}
\label{subsec:r91-threshold-identity}

The bound \(\beta_4=1/256\) does not occur once.  The source states in
Definition~\ref{def:v470-r13-local-relay-threshold}
\[
\Theta_v:=\beta_4^{r12}=\frac1{256},
\]
making the residual \(C4\) bound and the local ridge-front relay threshold the same quantity
by definition.  The detector chain carries the same value as \(\theta_{\rm det}=1/256\), and
r81 types it explicitly:
\[
\varepsilon_{\mathsf M}^{\rm len}=\frac{\ell_{\rm P}}{L_{\rm ref}}
\qquad\text{or, dimensionlessly,}\qquad
\theta_{\rm det}=\frac1{256},
\]
with the accompanying sentence \emph{``These are terminal scales, not exact all-readout
zeros.''}

Two readings are possible and should be distinguished.  Either \(1/256\) is one terminal
scale appearing in three roles --- seam bound, relay threshold, detector threshold --- or it is
three quantities that happen to share a dyadic value.  The source settles the first two by
definition; for the third it supplies equality of value and equality of role, but no explicit
identification.  The argument below needs only the first two plus the r81 typing.

\subsection*{r91 why spending is not available}
\label{subsec:r91-why-not-spendable}

\paragraph{Theoremlet r91.}
Spending \(\beta_4\) --- declaring the residual seam absorbed --- requires treating a residual
at the terminal scale as an exact zero.  Under the matter-frame typing of r81 that inference
is unavailable:
\[
\Phi\le\varepsilon_{\mathsf M}
\ \Longrightarrow\
0_{\mathsf M}^{\varepsilon}
\ \Longrightarrow\
0_{\rm termdata},
\qquad\text{but}\qquad
0_{\mathsf M}^{\varepsilon}\not\Rightarrow 0_{\rm all}.
\]
Hence \(Q_4\) is not blocked by a missing tensor or an unbuilt selector.  It is blocked by the
same typing that blocks \(C5=0\), and no refinement of the localization can remove that
block.

\paragraph{Proof.}
The residual seam is bounded by \(\beta_4\), and \(\beta_4=\Theta_v\) is a terminal threshold
of the declared chain.  A spending step asserts that the seam contributes nothing, that is
\(S_{4,\rm out}^{\rm collar}=0\) as an all-readout statement.  What the bound supports is
\(S_{4,\rm out}^{\rm collar}\le\beta_4\), which by
Theoremlet~r81 yields \(0_{\mathsf M}^{\varepsilon}\) and terminal data, and explicitly not
\(0_{\rm all}\).  The step from the bound to the assertion is exactly the forbidden
promotion. \(\square\)

\subsection*{r91 status reclassification}
\label{subsec:r91-status-reclassification}

Using the r86 vocabulary, the queue entry changes class without changing value:
\[
Q_4:\quad
\mathsf{inadmissible}_{\rm policy}
\ \longmapsto\
\mathsf{Excluded}_{\rm typed}.
\]
The practical difference matters for anyone planning further work.  Under the old reading,
\(Q_4\) was a branch one might open later, once a suitable comparison tensor had been built.
Under the corrected reading, building that tensor would not help: the obstruction is not the
absence of a channel but the typing of the scale at which the residual sits.  A future
\(C4\) closure, if one exists, must come from a route that does not pass through a
terminal-scale bound --- structurally the same requirement that r88 placed on shape (iii) for
\(C5\).

\subsection*{r91 physical reading}
\label{subsec:r91-physical-reading}

The formal result has a reading in the physical layer that the document otherwise keeps
separate, and it is recorded here as interpretation, not as a claim.

The commutator component was routed rather than eliminated.  In the compensator language of
the main line, where \(C=\tfrac12(\mathrm{id}+J)\) splits the space into symmetric and
antisymmetric parts, an antisymmetric transport remainder is not something a change of chart
removes; it is read through a comparison channel and kept.  That is the expected behaviour of
a noncommutative structure, and r07 treats it accordingly.

What survives as \(C4\) is a seam at the bounded--exterior interface, bounded by the terminal
resolution.  Read physically, the statement of r91 is that a residual sitting at the
resolution scale of the readout cannot be declared absent by that same readout.  This is the
same discipline that r81 expresses by pairing \(\theta_{\rm det}\) with \(\ell_{\rm P}/L_{\rm
ref}\): below the terminal scale the frame records a zero, and that record is a statement
about the frame, not about the object.

No claim is made that \(\beta_4\) equals a physical constant, that \(1/256\) has dimensional
meaning, or that the seam corresponds to a measurable quantity.

\subsection*{r91 ledger and guard update}
\label{subsec:r91-ledger-guard}

\(Q_4\) is executed and terminates negatively:
\[
R_{Q_4}^{r91}=(0,0,1)
\]
for (commutator component already routed, absorbability already settled, spending shown
unavailable).  The third component is recorded as \(1\) because a negative result closes the
branch, not the obstruction: \(C4\) remains open in the exact/global ledger.

The guard vector is unchanged, and pointedly so:
\[
R_{\rm glob}^{(r91,+)}=(0,1,1,1,1),
\qquad
\Gamma_{\rm guard}^{r91}=(1,1,1,1,1),
\qquad
\beta_4=\frac1{256}\ \text{unspent}.
\]
\(\beta_4\) remains unspent not as a reserve held back, but because the spending operation is
not available in the declared typing.  The layer does not imply exact finite \(C5\) zero
closure, exact all-readout terminality, \(C6/T2\), a global photon-clock/RH theorem, or
global all-readout closure.


% ============================================================================
\clearpage
\section*{v4.7.0-r92 Source-Only: Sharpened Reconstruction of \texorpdfstring{\(\Delta_\gamma<1/5\)}{Delta gamma < 1/5}}
\addcontentsline{toc}{section}{v4.7.0-r92 Source-Only: Sharpened Reconstruction of Delta gamma < 1/5}
\label{sec:r92-delta-gamma-sharpened}
% ============================================================================

\paragraph{Purpose.}
r90 reconstructed the bound \(\Delta_\gamma<1/5\) to within a rounding: the coherence
argument gave \(2/9=0.2222\ldots\), and \(1/5\) was read as a conservative rounding of it.
This layer replaces that with an exact account.  The constant is not a rounding.  It is the
relative gain reserve of a single named chart, and it can be traced through the arithmetic
that produces it.  As in r90 this is a reconstruction, not a recovery of the r65 text, which
a fourth search has again failed to locate.

\subsection*{r92 the arithmetic signature}
\label{subsec:r92-arithmetic-signature}

Every threshold in the declared chain is dyadic:
\[
\beta_4=\Theta_v=\theta_{\rm det}=\frac1{256}=2^{-8},
\qquad
a=\frac1{64}=2^{-6},
\qquad
\frac{\Theta_v}{a^2}=16=2^4 .
\]
The bound \(\Delta_\gamma<1/5\) is the only non-dyadic constant in the layer.  A factor of
five cannot arise from products and quotients of powers of two, so it must be produced by a
division that introduces it.  In the finite chart family there is exactly one such division.
The overlap flux of \(V3\) is
\[
S_{V3}=N+\kappa N(N-1)=5+\tfrac14\cdot5\cdot4=10,
\]
and the required gain is
\[
G_{V3}^{\rm req}=\frac{16}{S_{V3}}=\frac{16}{10}=\frac85 .
\]
The five enters here, by cancellation of \(16/10\), and nowhere else: \(V5\) shares
\(S=10\) and therefore the same \(8/5\), while every other chart yields a denominator free of
five.

\subsection*{r92 the perturbation picture}
\label{subsec:r92-perturbation-picture}

r90 let the perturbation act on the coherence factor \(K\) inside
\(\mathcal C_\gamma=1+\eta K\), giving \(\mathrm{Var}(\varphi)\le1-\rho_V\).  A second reading
is available and is the one the constant selects.  At full coherence
\(\mathcal C_\gamma=G_V^{\rm avail}\); a frequency spread lowers the probe value itself, and
the window survives while the probe stays above \(G_V^{\rm req}\).  The admissible
\emph{relative} drop is therefore
\[
\Delta_\gamma^{\max}(V)=1-\frac{G_V^{\rm req}}{G_V^{\rm avail}} .
\]

\paragraph{Theoremlet r92.}
For the charts carrying windows, the admissible relative drop is
\[
\begin{array}{c|c|c|c|c}
\text{chart} & \text{chain} & G^{\rm req} & G^{\rm avail} & 1-G^{\rm req}/G^{\rm avail}\\
\hline
V3 & \mathcal R_A & 8/5 & 2 & 1/5\\
V5 & \mathcal R_A & 8/5 & 5/2 & 9/25\\
V8 & \mathcal R_B & 16/9 & 2 & 1/9
\end{array}
\]
Minimized over the primary relay chain \(\mathcal R_A\) of
Section~\ref{sec:v470-r15-zero-clock-phase-window-relay-stitching}, this is
\[
\min_{V\in\mathcal R_A}\Bigl(1-\frac{G_V^{\rm req}}{G_V^{\rm avail}}\Bigr)
=\min\Bigl\{\frac15,\frac9{25}\Bigr\}
=\frac15 ,
\]
attained at \(V3\), which is the weakest assisted chart of that chain.

\paragraph{Proof.}
Substitute the r14 parameters into \(S=N+\kappa N(N-1)\) and \(G^{\rm req}=16/S\), then form
the quotient.  For \(V3\): \(S=10\), \(G^{\rm req}=8/5\), \(G^{\rm req}/G^{\rm avail}=4/5\),
so the admissible drop is \(1/5\).  For \(V5\): \(G^{\rm req}/G^{\rm avail}=16/25\), giving
\(9/25>1/5\).  The minimum over the two is \(1/5\). \(\square\)

\subsection*{r92 comparison with r90}
\label{subsec:r92-comparison}

The two reconstructions differ in where the perturbation is applied and in how well they fit:
\[
\begin{array}{l|c|c}
& \text{value on }\mathcal R_A & \text{fit to }1/5\\
\hline
\text{r90, spread acts on coherence }K & 2/5 & \text{no}\\
\text{r90, minimum over all three charts} & 2/9 & \text{by rounding}\\
\text{r92, spread lowers the probe relatively} & 1/5 & \text{exact}
\end{array}
\]
An exact hit through a chain of four elementary steps --- packet count, overlap flux, required
gain, relative reserve --- is considerably stronger evidence than agreement after rounding.
r90 is therefore superseded as an account of the constant, though its coherence bound remains
correct as a statement in its own right.

\subsection*{r92 what remains uncertain}
\label{subsec:r92-uncertain}

Two points are not settled by the fit.

First, the minimum is taken over \(\mathcal R_A\) alone.  Including \(V8\), which belongs to
the separate chain \(\mathcal R_B\) in the r15 stitching, would give \(1/9\) instead.  The
restriction is consistent with the guard being written \(W_{\rm clk}^{\rm loc}\), a local
statement, and with the two chains being separate at the time the windows were extracted;
but the source does not state it.

Second, the linear response of \(\mathcal C_\gamma\) to the frequency spread is assumed, not
derived.  What the fit establishes is that the constant \(1/5\) is the relative gain reserve
of \(V3\); it does not establish that the original argument reached it by this route.

\subsection*{r92 ledger and guard update}
\label{subsec:r92-ledger-guard}

The reconstruction improves but does not close the provenance entry:
\[
R_{\rm rec}^{r92}=(0,0,1)
\]
for (arithmetic origin of the five located, exact value derived, r65 text still absent).  The
third component can only be closed by the original text.

r92 spends no guard:
\[
R_{\rm glob}^{(r92,+)}=(0,1,1,1,1),
\qquad
\Gamma_{\rm guard}^{r92}=(1,1,1,1,1).
\]
All statements are finite, local and chart-bound.  The layer makes no spectral identification
between Riemann zeros and photon ticks and does not imply exact finite \(C5\) zero closure,
exact all-readout terminality, \(C6/T2\), a global photon-clock/RH theorem, or global
all-readout closure.


% ============================================================================
\clearpage
\section*{v4.7.0-r93 Source-Only: The Two-Sided Typing Read as a Routing Grammar}
\addcontentsline{toc}{section}{v4.7.0-r93 Source-Only: The Two-Sided Typing Read as a Routing Grammar}
\label{sec:r93-routing-reading}
% ============================================================================

\paragraph{Purpose.}
The two-sided structural tensor notation of
Definition~\ref{def:two-sided-structural-tensor-interface} has been in continuous use since
r26 without being named for what its index slots actually carry.  Read from the definition
alone, the front indices look like role annotations attached to an otherwise ordinary tensor
display.  Read from the sites that use them, they are addresses in a labelled graph, and the
contraction discipline is the forwarding rule of that graph.  This layer records the reading.
It introduces no object, no new contraction and no new claim; it names a structure already
present in the source since r26.

\subsection*{r93 what the index slots carry}
\label{subsec:r93-slot-content}

The decisive evidence is in the application layer, not in the definition layer.  The r26
transition typing (\S\ref{subsec:v470-r26-two-sided-transition-typing}) types four atlas
cells as
\[
{}^{\partial\mathrm{src}}{}_{\mathrm{loc}}\mathcal A_0{}^{V0}_{\mathrm{relay}},
\quad
{}^{\mathrm{core}}{}_{\mathrm{loc}}\mathcal A_1{}^{V5}_{\mathrm{relay}},
\quad
{}^{\mathrm{shadow}}{}_{\mathrm{loc}}\mathcal A_2{}^{V7}_{\mathrm{bridge}},
\quad
{}^{\gamma\mathrm{clock}}{}_{\mathrm{loc}}\mathcal A_3{}^{V9}_{\mathrm{tail}} .
\]
The back upper slot carries a chart name from the finite scan of
\S\ref{sec:v470-r14-finite-valley-front-chart-scan}; the back lower slot carries an edge
kind, one of \(\mathrm{relay}\), \(\mathrm{bridge}\), \(\mathrm{tail}\).  The accompanying
r26 audit ledger lists the transitions \(\tau_{01},\tau_{12},\tau_{23}\) by source type,
target type and status.  Node, edge kind, source, target, status: this is the data of a
labelled directed graph, recorded in index position.

The r50 refinement (\S\ref{subsec:r50-two-sided-domain-typing}) makes the four slots explicit
and names them:
\[
{}^{L_x}{}_{G_x}\mathcal O_x{}^{R_x}_{M_x},
\]
with \(L_x\) the local readout/source side, \(G_x\) the guard or selector side, \(R_x\) the
relay/structure side, and \(M_x\) the monitor codomain side.  The r49 recap
(\S\ref{subsec:r49-two-sided-typing-recap}) shows the same grammar carrying a transport
condition: a local boundary source may feed a local relay monitor only after the guard
projection has separated it from \(C4\), \(C6/T2\) and the all-readout channel.

\paragraph{Theoremlet r93.}
The two-sided display grammar of
Definitions~\ref{def:two-sided-structural-tensor-interface}
and~\ref{def:two-sided-admissible-contractions}, as instantiated from r26 onward, is an
interface grammar over a labelled directed graph.  Under the correspondence
\[
\begin{array}{lcl}
\text{back upper slot} &\longleftrightarrow& \text{node of the finite chart family}\\
\text{back lower slot} &\longleftrightarrow& \text{edge kind at that node}\\
\text{front upper slot} &\longleftrightarrow& \text{local source/readout attachment}\\
\text{front lower slot} &\longleftrightarrow& \text{guard or selector governing the attachment}
\end{array}
\]
the admissible-contraction rules of
Definition~\ref{def:two-sided-admissible-contractions} read as a forwarding discipline:
classical upper--lower contraction is transport along a declared edge, and the prohibition on
same-variance pairing without a declared comparison tensor or selector is the prohibition on
joining two channels without a declared admission element.

\emph{Proof sketch.}
The correspondence is read off the instantiations rather than postulated.  In every
two-sided display occurring from r26 onward the back upper slot is filled by a member of the
r14 chart family or by an obstruction name carried in the same position, and the back lower
slot is filled by an edge kind or monitor role.  The r26 transition ledger supplies source,
target and status for each transition, which is the incidence data of the graph.  The
forwarding reading of the contraction rules is a restatement, not an extension: rule~1 of
Definition~\ref{def:two-sided-admissible-contractions} permits exactly the pairings in which
one slot is a declared codomain of the other, and rules~2--4 admit a pairing only in the
presence of a named \(G_{\alpha\beta}\), \(G^{\alpha\beta}\) or \(\chi_{mn}\).  No pairing is
added and none is removed.  \(\square\)

\subsection*{r93 consequence for the relay arithmetic}
\label{subsec:r93-admission-reading}

Under the routing reading the finite relay layer of
\S\ref{sec:v470-r15-zero-clock-phase-window-relay-stitching} is an admission calculation.
With the overlap flux and required gain
\[
S_V=N+\kappa N(N-1),
\qquad
G_V^{\rm req}=\frac{16}{S_V},
\]
a chart is admissible exactly when its available gain covers its required gain, and the
window criterion \(16/G_{\rm avail}\le S_V<16\) is the corresponding capacity band.  The
three charts recorded as blocked in the r14 scan are those whose requirement exceeds what the
node supplies.  The guarded micro-lift of
\S\ref{sec:v470-r17-atlas-gap-bridge-v9-microlift}, which carried \(V9\) from \((2,1/2,3)\) to
\((3,1/2,3)\) and obtained \(H_9'=18>16\), is a capacity increase on a single node, executed
until the link carries and no further.

This also sharpens the reading of \S\ref{sec:r92-delta-gamma-sharpened}.  The relative reserve
\[
1-\frac{G_{V}^{\rm req}}{G_{V}^{\rm avail}},
\qquad
\text{attained as } 1-\tfrac45=\tfrac15 \text{ at } V3,
\]
is the margin by which the weakest assisted node exceeds its own requirement.  Minimising it
over the chain is the ordinary statement that a chain is limited by its weakest link.  The
r92 derivation is unchanged; only its interpretation is now stated in the vocabulary the
notation already uses.

\subsection*{r93 conservativity is not inertness}
\label{subsec:r93-conservativity}

Proposition~\ref{prop:two-sided-conservative-expansion} states that every well-formed
two-sided display expands into the inherited QAM-over-ZF language and that the notation cannot
by itself introduce a new primitive tensor kind, a new admissibility relation or a new
theorematic certificate.  That proposition is not weakened here, and this layer relies on it.

It is worth recording that conservativity in that sense does not mean the notation is
decorative.  A routing table is likewise definable in the language of the objects it routes
between, and it still determines what reaches what.  The two statements are compatible: the
graph structure carried in the index slots adds no mathematical primitive, and it does
determine which transports are admissible.  Reading conservativity as inertness is a misreading
that this layer is written to prevent.

\subsection*{r93 what is not claimed}
\label{subsec:r93-not-claimed}

The front slots carry readout, frame, guard and selector data.  They are not asserted to
transform tensorially under any coordinate change, and no transformation law is supplied or
implied for them.  The graph reading is a reading of the declared typing; it is not a claim
that the chart family is a physical network, that the edges are physical channels, or that
the node capacities are measurable quantities.  The notation is used locally and sparingly in
the present release line, as recorded at its introduction, and r93 does not rewrite the
monolith into it.

\subsection*{r93 ledger and guard update}
\label{subsec:r93-ledger-guard}

The naming is complete for the sites that instantiate the notation; the source is not
uniform in its use:
\[
R_{\rm route}^{r93}=(0,0,1)
\]
for (slot semantics identified, forwarding reading of the contraction rules stated, uniform
use across the monolith not performed).  The third component is a presentation entry and
carries no obstruction.

r93 spends no guard:
\[
R_{\rm glob}^{(r93,+)}=(0,1,1,1,1),
\qquad
\Gamma_{\rm guard}^{r93}=(1,1,1,1,1).
\]
All statements are finite, local and chart-bound.  The layer makes no spectral identification
between Riemann zeros and photon ticks and does not imply exact finite \(C5\) zero closure,
exact all-readout terminality, \(C6/T2\), a global photon-clock/RH theorem, or global
all-readout closure.

% ============================================================================
\clearpage
\section*{v4.7.0-r94 Source-Only: Provenance of the C5 Contraction Factor 3/4}
\addcontentsline{toc}{section}{v4.7.0-r94 Source-Only: Provenance of the C5 Contraction Factor 3/4}
\label{sec:r94-c5-factor-provenance}
% ============================================================================

\paragraph{Purpose.}
Since r09 the bounded C5 collar runs on the update
\(\mathcal U_5^{\rm bd}:\Phi\mapsto\tfrac34\Phi\), and r86 inherits it without revisiting
where the factor comes from.  This layer traces the factor to its first occurrence, records
what the source does and does not supply as its justification, and separates the results that
depend on the exact value \(3/4\) from those that hold for every \(q\in(0,1)\).  The question
was raised in an external OP5/QAM working note kept in the repository folder
\texttt{analysis/} (sections ``Re-audit of the origin and necessity of the factor three
quarters'' and ``Step twenty does not identify a unique coefficient'' of
\texttt{OP5\_QAM\_Residual\_Lift.tex}), which was written against the RC19 source and did not
see r85--r93.  r94 checks its two findings against the current monolith and records them in
source.  It introduces no new object and no new claim about \(C5\).

\noindent\textbf{Source-only status.}
This r94 layer is source-only.  It does not build a PDF, does not assign a new DOI,
and does not authorize a public release.  The intended reserved v4.7.0 Version DOI remains
\href{https://doi.org/10.5281/zenodo.20232911}{10.5281/zenodo.20232911}; the previous
baseline DOI \href{https://doi.org/10.5281/zenodo.20117361}{10.5281/zenodo.20117361}
remains provenance only.

\subsection*{r94 provenance trace}
\label{subsec:r94-provenance-trace}

The C5 chain introduces and propagates the factor as follows.
\[
\begin{array}{l|l|l|l}
\text{layer} & \text{object} & \text{value} & \text{how the value enters}\\
\hline
\text{r04} & \text{Def.~\ref{def:v470-r04-finite-c5-witness-ledger}, row }W_5^{\rm bd}
 & \Phi_5=\tfrac34,\ \Delta_5=-\tfrac14 & \text{ledger entry}\\
\text{r05} & \text{Def.~\ref{def:v470-r05-finite-c5-iteration-ledger}, row }Z_5^{\rm bd2}
 & \Phi_5^{(2)}=\tfrac9{16} & \text{ledger entry}\\
\text{r05} & \text{same ledger, row }Z_5^{\rm collar}
 & \Phi_5^{(2)}=\tfrac58 & \text{ledger entry, ratio }\tfrac56\\
\text{r06} & \text{Def.~\ref{def:v470-r06-c5-decision-alphabet}}
 & q_1=q_2\le\tfrac34 & \text{decision criterion}\\
\text{r07} & \text{Def.~\ref{def:v470-r07-third-step-c5-probe}}
 & \Phi_5^{(3)}=\tfrac{27}{64} & \text{probe entry}\\
\text{r08} & \text{Lem.~\ref{lem:v470-r08-fourth-c5-bounded-step}}
 & \Phi_5^{(4)}=\tfrac{81}{256} & \text{computed by the candidate update}\\
\text{r09} & \text{Def.~\ref{def:v470-r09-bounded-c5-operator-domain}}
 & \mathcal U_5^{\rm bd}=\tfrac34\,\cdot & \text{operator definition}
\end{array}
\]
The r04 value \(\Delta_5(W_5^{\rm bd})=-\tfrac14\) is defined as the difference of two listed
scores; it is not an independent computation of a removed quarter.

\paragraph{Theoremlet r94.1 (provenance).}
In the present source the C5 contraction factor \(3/4\) is a declared value, not a derived
one.  Specifically:
\begin{enumerate}[leftmargin=2em]
 \item its first occurrence in the C5 sense is the r04 ledger entry \(\Phi_5(W_5^{\rm bd})=3/4\);
 \item the trace values \(3/4\), \(9/16\), \(27/64\) of r04, r05 and r07 are ledger or probe
 entries, and the r08 value \(81/256\) is produced by applying the candidate update itself
 (Proposition~\ref{prop:v470-r08-finite-induction-candidate-status} correctly classifies it as a
 prediction, \(\mathsf{cand}_{\rm ind}\));
 \item no layer supplies a mechanism --- a projection identity, an energy partition, a dynamical
 rate --- that produces the value \(3/4\) from MM/HC, collar or readout data.
\end{enumerate}
Consequently the phrase ``four independent bounded lowering steps'' in the proof of
Lemma~\ref{lem:v470-r09-c5-finite-stage-admissibility} is to be read as ``four consecutive
bounded lowering steps'': three of them are recorded entries and the fourth is the operator's
own prediction.  The lemma itself is unaffected, because its proof of admissibility rests on
positivity and domain preservation, which hold for the declared factor regardless of its
origin.

\paragraph{Proof sketch.}
Items 1 and 2 are read off the cited definitions and lemmas.  For item 3 the working file was
searched in full for the forms \texttt{\string\tfrac34}, \texttt{\string\frac34},
\texttt{\string\frac\{3\}\{4\}}, \texttt{3/4} and \texttt{quarter}.  The occurrences before r04
belong to other objects: the quarter-turn \(R_W\) of the minimal two-directional leakage plane
(Corollary~\ref{cor:working-first-place-where-quadrature-becomes-possible}), which is an
isometry and therefore cannot contract a norm, and the equal quarter-load shells of
Corollary~\ref{cor:op1-quarter-shells}, which concern OP~1 and are explicitly declared not
necessary.  From r05 onward every C5 occurrence either restates the r04--r09 trace or applies
Theorem~\ref{thm:v470-r09-c5-bounded-operator-theorem}; the r86 inventory
(\S\ref{subsec:r86-collar-inventory}) lists Definition~\ref{def:v470-r09-bounded-c5-operator-domain}
as the origin of \(\mathcal U_5^{\rm bd}\) and states that everything in it is inherited.  The
only occurrences that use the word ``calibrated'' concern the relay coupling \(\kappa=3/4\) of
the comfort-calibrated clock seed \((N,\kappa,G)=(5,3/4,1)\), a different parameter, chosen with
positive margin above the requirement \(\kappa\ge11/20\); no line of the source identifies
\(\kappa\) with the C5 ratio.  \(\square\)

\paragraph{Remark (why a single projection does not explain the factor).}
The working note observes that an orthogonal least-squares correction \(r'=r-d\) with
\(\langle r-d,d\rangle=0\) gives \(\|r'\|^2=\|r\|^2-\|d\|^2\), hence the score ratio \(3/4\) when
the removed component carries one quarter of the residual energy.  This explains one step at
most: the projection is idempotent, so applying it again changes nothing, and a geometric
recurrence would need a refreshed fit direction with a new one-quarter identity at every stage.
The source declares neither.  The remark is recorded as a constraint on any future derivation,
not as a candidate mechanism.

\subsection*{r94 what the detector stage determines}
\label{subsec:r94-detector-stage-interval}

r86 computes the detector stage from the factor
(\S\ref{subsec:r86-threshold-function}).  The converse direction is not recorded there.

\paragraph{Theoremlet r94.2.}
Let \(\Phi^{(n)}=q^n\) with \(q\in(0,1)\) and detector threshold \(\theta=1/256\).  The first
stage with \(\Phi^{(n)}\le\theta\) equals \(20\) if and only if
\[
\Bigl(\tfrac1{256}\Bigr)^{1/19}<q\le\Bigl(\tfrac1{256}\Bigr)^{1/20},
\qquad\text{numerically}\qquad
0.746879\ldots<q\le0.757858\ldots .
\]
Hence the recorded detector stage \(n=20\) is consistent with \(q=3/4\) but does not single it
out: every factor in a half-open interval of width \(0.01098\ldots\) produces the same stage.

\paragraph{Proof sketch.}
For \(q\in(0,1)\) the sequence \(q^n\) is strictly decreasing, so the first crossing is at
stage \(20\) exactly when \(q^{20}\le1/256<q^{19}\).  Taking the \(19\)th and \(20\)th roots gives
the interval.  For \(q=3/4\) the exact rational check is
\((3/4)^{19}=0.004228\ldots>1/256\ge0.003171\ldots=(3/4)^{20}\), as in the remark to
Theoremlet r86.1.  \(\square\)

\paragraph{Illustration from the r05 ledger.}
The tolerated collar row \(Z_5^{\rm collar}\) of
Definition~\ref{def:v470-r05-finite-c5-iteration-ledger} has ratio \((5/8)/(3/4)=5/6\).  Iterated
at that ratio the score would first reach \(1/256\) at stage \(31\), not \(20\).  The stage
\(n=20\) therefore records the choice of the strict row \(Z_5^{\rm bd2}\) over the tolerated row,
which r05 makes by minimality, and not a property of the collar as such.  The figures \(20\)
and \(31\) were checked by iterating in exact rational arithmetic.

\subsection*{r94 which results depend on the exact value}
\label{subsec:r94-value-dependence}

\[
\begin{array}{l|l}
\text{holds for every } q\in(0,1) & \text{depends on } q=3/4 \text{ or the r94.2 interval}\\
\hline
\text{finite positivity, Lem.~\ref{lem:v470-r11-finite-positive-guard}}
 & \text{stage } n=20 \text{ at } \theta=1/256 \text{ (interval)}\\
\text{limit zero, Thm.~\ref{thm:v470-r09-c5-bounded-operator-theorem}}
 & \text{halving cost } \ln2/\ln(4/3) \text{ of r86.1}\\
\text{finite-zero exclusion, Theoremlet r86.2}
 & \text{instances } c_\star(1),\dots,c_\star(5) \text{ of r88.1}\\
\text{subtractive obstruction, Theoremlet r88.2} & \\
\text{r88 shape (iii): } (\alpha),(\beta),(\delta) &
 \text{r88 shape (iii): } (\gamma) \text{ (interval)}
\end{array}
\]
r86.2 already notes that the exclusion depends only on \(q>0\).  r94 extends that observation to
the whole status picture: none of the structural conclusions of r86 and r88 uses the particular
value, and the one guarded claim that mentions it,
\(C5_{\rm det/eff}\) at \(\theta=1/256\) after \(n=20\), needs only the interval of
Theoremlet r94.2.

This is the same situation r88 records for the subtractive channel
(\S\ref{subsec:r88-free-parameter-trap}): a value that is fixed by declaration rather than by an
independent structural rule carries the content of the declaration, not of a derivation.  For
\(q\) this is harmless, because no guarded conclusion depends on more than the interval; it
becomes relevant only if a later layer wants to treat \(3/4\) itself as a physical or
geometric constant.

\subsection*{r94 ledger and guard update}
\label{subsec:r94-ledger-guard}

\[
R_{\rm prov}^{r94}=(1,0,0)
\]
for (independent derivation of the factor \(3/4\) not available, detector-stage interval
determined, value-dependence of the C5 and r86/r88 results sorted).  The first component is
an open provenance item.  It is not an obstruction, it does not enter the r48 branch queue, and
none of the guarded claims depends on closing it.

r94 spends no guard and changes no status:
\[
R_{\rm glob}^{(r94,+)}=(0,1,1,1,1),
\qquad
\Gamma_{\rm guard}^{r94}=(1,1,1,1,1),
\qquad
\beta_4=1/256\ \text{unspent}.
\]
The admissible claim \(C5_{\rm det/eff}\) at \(\theta=1/256\) after \(n=20\) is unchanged.  The
layer does not imply exact finite \(C5\) zero closure, exact all-readout terminality,
\(C6/T2\), a global photon-clock/RH theorem, or global all-readout closure.

% ============================================================================
\clearpage
\section*{v4.8.0-r01 Source-Only: Planck-Readout Metric Layer}
\addcontentsline{toc}{section}{v4.8.0-r01 Source-Only: Planck-Readout Metric Layer}
\label{sec:v480-r01-planck-readout-metric-layer}
% ============================================================================

\paragraph{Purpose.}
This layer opens the v4.8.0 line.  It integrates the Planck-readout candidate prepared on the
repository research branch for v4.8.0-r01 and states it in the frame-typed vocabulary of
r81--r82.  The candidate separates the physical Planck ruler from the scale carried by a
lower interface or transport layer: what is invariant is not the number \(\lP\) in every
layer, but its guarded recovery under an admissible readout.  The monolith already contains
one instance of this pattern, the photon-frame Planck length of
Theorem~\ref{thm:planck-photon}; r01 names the general form and fixes how its sign is to be
read relative to the r82 signature convention.  The field-equation refactor announced in the
working title of the research line is not carried out here.

\noindent\textbf{Source-only status.}
This v4.8.0-r01 layer is source-only.  It does not build a PDF, does not assign a new DOI,
and does not authorize a public release.  The v4.7.0 Version DOI
\href{https://doi.org/10.5281/zenodo.20232911}{10.5281/zenodo.20232911} and the Concept DOI
\href{https://doi.org/10.5281/zenodo.19210728}{10.5281/zenodo.19210728} are unchanged by it.

\subsection*{v4.8.0-r01 readout-calibrated Planck scale}
\label{subsec:v480-r01-calibrated-scale}

\paragraph{Definition v4.8.0-r01.1.}
Let \(R\) be an admissible readout or interface mode and let \(\lambda_R\in\mathbb R\setminus\{0\}\)
be the local calibration factor induced by the MM/HC/QAM transport from the structural
interface to the physical readout \(R\).  Write
\[
\ell_{\rm eff}(R):=\lP\,\lambda_R,
\qquad
o_R:=\operatorname{sgn}\lambda_R\in\{+1,-1\},
\qquad
|\ell_{\rm eff}(R)|=\lP\,|\lambda_R| .
\]
\(\ell_{\rm eff}(R)\) is a signed interface coordinate.  Its magnitude is the length scale the
readout \(R\) resolves; its sign \(o_R\) is an orientation code.

\paragraph{Theoremlet v4.8.0-r01.}
Under Definition v4.8.0-r01.1 the following hold.
\begin{enumerate}[leftmargin=2em]
 \item \emph{Matter recovery.}  If \(R_{\mathsf M}\) is a stable matter readout with
 \(\lambda_{R_{\mathsf M}}=1+\varepsilon_{\mathsf M}\) and
 \(|\varepsilon_{\mathsf M}|\le\delta_{\mathsf M}<1\), then
 \(o_{R_{\mathsf M}}=+1\) and
 \(\bigl|\ell_{\rm eff}(R_{\mathsf M})-\lP\bigr|\le\delta_{\mathsf M}\lP\): the ordinary Planck
 scale is recovered up to the declared local calibration error.
 \item \emph{Photon instance.}  The photon-frame Planck length of
 Theorem~\ref{thm:planck-photon} is the instance
 \(\lambda_{\gamma}=\gL^{-3/2}\), \(o_\gamma=+1\), of the same definition.
 \item \emph{Opposite orientation.}  If \(R_\tau\) is an opposite-orientation readout with
 \(\lambda_{R_\tau}=-\mu_{R_\tau}\), \(\mu_{R_\tau}>0\), then
 \(\ell_{\rm eff}(R_\tau)=-\lP\mu_{R_\tau}\) while
 \(|\ell_{\rm eff}(R_\tau)|=\lP\mu_{R_\tau}>0\).  The negative sign is the orientation code
 \(o_{R_\tau}=-1\) and not a negative physical distance:
 \[
 o_{R}=-1\ \not\Rightarrow\ L_{\mathsf M}<0 .
 \]
\end{enumerate}

\paragraph{Proof sketch.}
Item 1: \(\delta_{\mathsf M}<1\) gives \(\lambda_{R_{\mathsf M}}>0\), and
\(\ell_{\rm eff}(R_{\mathsf M})-\lP=\lP\varepsilon_{\mathsf M}\).  Item 2: substituting
\(\lambda_\gamma=\gL^{-3/2}\) into the definition reproduces the display of
Theorem~\ref{thm:planck-photon}; nothing is proved anew and the status of that theorem is
inherited unchanged.  Item 3: the definition separates the signed coordinate from its
magnitude, and the magnitude is positive whenever \(\mu_{R_\tau}>0\).  The non-implication is
the r82 guard \(\ell_{\mathsf T}^2<0\not\Rightarrow L_{\mathsf M}<0\) restated for the
orientation code.  \(\square\)

\subsection*{v4.8.0-r01 orientation code and r82 signature are different data}
\label{subsec:v480-r01-orientation-vs-signature}

r82 types the tachyon-frame scale through the sign of the \emph{interval},
\(\ell_{\mathsf T}^2=-\lambda_{\mathsf T}^2\) with \(\lambda_{\mathsf T}>0\).  r01 types an
opposite-orientation readout through the sign of a \emph{linear} coordinate,
\(o_{R_\tau}=-1\).  The two are not the same bit: squaring the r01 coordinate gives
\(\ell_{\rm eff}(R_\tau)^2=\lP^2\mu_{R_\tau}^2>0\), which is not the r82 interval.  This layer
therefore does not identify them.

If a later layer calibrates the tachyon-frame native magnitude as
\(\lambda_{\mathsf T}=\lP\mu_{R_\tau}\), the readout \(R_\tau\) carries two independent records,
\[
\bigl(o_{R_\tau},\ \operatorname{sgn}\ell_{\mathsf T}^2\bigr)=(-1,-1),
\qquad
\text{magnitude } \lP\mu_{R_\tau},
\]
an orientation code from r01 and a signature marker from r82.  Neither record is a negative
material length, and neither is derivable from the other without a declared comparison map.
This is the r93 discipline: two channels are not joined without an admission element.

\paragraph{Remark (two-sided notation).}
The research note records the calibration schematically as
\({}^{R}{}_{\mathrm{phys}}\mathcal T^{\mathrm{struct}}{}_{\mathrm{HC/MM}}\).  Under the slot
reading of r93 the front lower slot carries a guard or selector, so the matching r93 form is
\({}^{R}{}_{G_R}\mathcal T^{\mathrm{struct}}{}_{\mathrm{HC/MM}}\) with \(G_R\) the calibration
guard \(|\varepsilon_R|\le\delta_R\).  The back slots here name a layer and a transport, not a
chart of the r14 family; the display is schematic and is not an instance of the r93 routing
correspondence.

\subsection*{v4.8.0-r01 what is not claimed}
\label{subsec:v480-r01-not-claimed}

The layer does not derive any \(\lambda_R\) except by citing
Theorem~\ref{thm:planck-photon}; for every other readout \(\lambda_R\) is a declared
calibration datum.  It does not assert that \(\lP\) is non-invariant as a physical constant,
does not promote a negative sign to an observable negative length, does not prove exact
through-collapse transport, and does not refactor the MM/HC field equations.

\subsection*{v4.8.0-r01 ledger and guard update}
\label{subsec:v480-r01-ledger-guard}

\[
R_{\rm Planck}^{v480\text{-}r01}=(0,0,1)
\]
for (calibrated scale defined with matter recovery and photon instance, orientation code
separated from the r82 signature, MM/HC field-equation refactor not performed).  The third
component is open research, not an obstruction.

The layer spends no guard:
\[
R_{\rm glob}^{(v480\text{-}r01,+)}=(0,1,1,1,1),
\qquad
\Gamma_{\rm guard}^{v480\text{-}r01}=(1,1,1,1,1),
\qquad
\beta_4=1/256\ \text{unspent}.
\]
It does not imply exact finite \(C5\) zero closure, exact all-readout terminality,
\(C6/T2\), a global photon-clock/RH theorem, or global all-readout closure.

% ============================================================================
\clearpage
\section*{v4.8.0-r02 Release Integration and Scope Refresh}
\addcontentsline{toc}{section}{v4.8.0-r02 Release Integration and Scope Refresh}
\label{sec:v480-r02-release-integration}
% ============================================================================

\paragraph{Purpose.}
This layer integrates the post-v4.7.0 source-only layers and the opening layer of the
v4.8.0 line into the v4.8.0 release, in the manner in which r84 integrated r81--r83 into the
v4.7.0 release scope.  It adds no mathematical content.

\noindent\textbf{Release status.}
Unlike the preceding layers, r02 is a release layer.  It accompanies the v4.8.0 Version DOI
\href{https://doi.org/10.5281/zenodo.23365244}{10.5281/zenodo.23365244} under the Concept DOI
\href{https://doi.org/10.5281/zenodo.19210728}{10.5281/zenodo.19210728}.  The v4.7.0 Version
DOI \href{https://doi.org/10.5281/zenodo.20232911}{10.5281/zenodo.20232911} becomes provenance.

\subsection*{v4.8.0-r02 integrated layers}
\label{subsec:v480-r02-integrated-layers}

\[
\begin{array}{l|l|l}
\text{layer} & \text{content} & \text{status change}\\
\hline
\text{r85} & \text{citation repair, unity-theorem scope separation} & \text{none}\\
\text{r86} & \text{C5 collar bookkeeping, finite-zero exclusion for } q>0 & \text{none}\\
\text{r87, r89} & \text{selector-window calibration and its correction} & \text{none}\\
\text{r88} & \text{non-collar routes: shapes (i), (ii) excluded, (iii) open} & \text{none}\\
\text{r90, r92} & \text{reconstruction of } \Delta_\gamma<1/5 & \text{none}\\
\text{r91} & \text{Q4 executed: } \beta_4 \text{ not spendable} & \text{none}\\
\text{r93} & \text{two-sided typing read as a routing grammar} & \text{none}\\
\text{r94} & \text{provenance of the C5 factor } 3/4 & \text{none}\\
\text{v4.8.0-r01} & \text{readout-calibrated Planck scale} & \text{none}
\end{array}
\]

The open items carried into v4.8.0 are: shape (iii) of r88, the independent derivation of
the C5 factor (r94), uniform use of the two-sided notation (r93), and the MM/HC field-equation
refactor (v4.8.0-r01).  None of them is an obstruction to a guarded claim.

\subsection*{v4.8.0-r02 guard update}
\label{subsec:v480-r02-guard-update}

The release spends no guard:
\[
R_{\rm glob}^{(v480\text{-}r02,+)}=(0,1,1,1,1),
\qquad
\Gamma_{\rm guard}^{v480\text{-}r02}=(1,1,1,1,1),
\qquad
\beta_4=1/256\ \text{unspent}.
\]
The safe release formulation remains: guarded effective/local and frame-typed closures, not
exact global closure.  The release does not imply exact finite \(C5\) zero closure, exact
all-readout terminality, \(C6/T2\), a global photon-clock/RH theorem, global all-readout
closure, or a negative physical length.
